\subsection{Options}\label{sec:options}
An option is a derivative that gives its holder the right, but not the obligation, to buy or sell an underlying asset at a fixed price. Options are traded on exchanges and also directly between two parties, which is referred to as over-the-counter trading. A \textbf{call option} is a right to buy the underlying asset and a \textbf{put option} is a right to sell it. The fixed price is called the strike price $K$, and the date on which the right expires is the expiry date. 

American options can be exercised at any point up to the expiry date. European options can only be exercised on the expiry date \parencite[31]{hullOptionsFuturesOther2022}. TTF gas options are European.

The payoff of an option is the value of the right at expiry. For a call option, the payoff is $\max(F_T - K, 0)$, where $F_T$ is the settlement price at expiry. The right is exercised only if $F_T > K$, otherwise the option is left unused.
For a put option, the payoff is $\max(K - F_T, 0)$. The payoff of an option is never negative \parencite[231]{hullOptionsFuturesOther2022}. When the option is bought, the buyer pays the seller an amount called the premium. The premium is the price of an option.

The relation between the current price and $K$ is described by three terms. A call option is \textbf{in the money} if $F_t > K$, \textbf{at the money} if $F_t = K$ and \textbf{out of the money} if $F_t < K$, where $F_t$ is the settlement price on trading day $t$. For a put option, the definitions are reversed \parencite[234]{hullOptionsFuturesOther2022}.

The settlement price at expiry is unknown when the option is bought. A risk-neutral world is a world in which investors do not require a higher expected return for taking on more risk. In a risk-neutral world, the expected return on any investment is the risk-free rate. Pricing an option under this assumption gives the correct price in the real world as well \parencite[292]{hullOptionsFuturesOther2022}. The price of a European call is its expected payoff in a risk-neutral world, discounted at the risk-free rate

\begin{equation}
    c = e^{-r T} \hat{\mathbb{E}}[\max(F_T - K, 0)],
    \label{eq:price-call}
\end{equation}

and the price of a European put is 

\begin{equation}
    p = e^{-r T} \hat{\mathbb{E}}[\max(K - F_T, 0)]
    \label{eq:price-put}
\end{equation}

where $\hat{\mathbb{E}}$ denotes the expected value in a risk-neutral world, $r$ is the risk-free rate and $T$ is the time to maturity. The option is valued at time zero, so that the expiry date and the time to maturity are both equal to $T$. The probabilities entering $\hat{\mathbb{E}}$ are those of a risk-neutral world and not those observed in the price series. The price of an option depends on the distribution of $F_T$, rather than on its expected value alone.

An option that is in the money would pay a positive amount at the current price, while an option that is out of the money would pay nothing. As a result, the value of an option depends entirely on the probability that the price moves past the strike before expiry.